\documentclass[11pt,a4paper]{article}
\usepackage[T1]{fontenc}
\usepackage[utf8]{inputenc}
\usepackage[brazilian]{babel}
\usepackage[default,type1,semibold]{inter}
\usepackage[a4paper,top=18mm,bottom=24mm,left=18mm,right=18mm,footskip=12mm]{geometry}
\usepackage[table]{xcolor}
\usepackage{amssymb,textcomp,array,longtable,enumitem,fancyhdr,lastpage,needspace,fontawesome5}
\usepackage[skins,breakable,raster]{tcolorbox}
\usepackage[hidelinks]{hyperref}
\usetikzlibrary{positioning,calc,babel}
\DeclareUnicodeCharacter{2192}{\ensuremath{\rightarrow}}
\DeclareUnicodeCharacter{2212}{\ensuremath{-}}
\DeclareUnicodeCharacter{2265}{\ensuremath{\geq}}
\DeclareUnicodeCharacter{2248}{\ensuremath{\approx}}
\DeclareUnicodeCharacter{03C3}{\ensuremath{\sigma}}

\definecolor{cor}{HTML}{B07D24}
\definecolor{suave}{HTML}{FAF1DF}
\definecolor{texto}{HTML}{0F172A}
\definecolor{cinza}{HTML}{5B6B84}
\definecolor{linha}{HTML}{D3DAE6}
\color{texto}
\setlength{\parindent}{0pt}
\setlength{\parskip}{4pt}
\linespread{1.1}
\setlength{\LTpre}{6pt}
\setlength{\LTpost}{10pt}
\renewcommand{\arraystretch}{1.45}
\arrayrulecolor{linha}
\setlist{itemsep=4pt,topsep=4pt,parsep=0pt}
\setlist[itemize]{label={\color{cor}\small\faCaretRight},leftmargin=1.4em}
% Passos numerados: número dentro de um círculo da cor da equipe.
\DeclareRobustCommand{\bolinha}[1]{\tikz[baseline=(n.base)]\node[circle,fill=cor,text=white,inner sep=1.6pt,minimum size=17pt,font=\footnotesize\bfseries](n){#1};}
\setlist[enumerate]{label=\bolinha{\arabic*},leftmargin=2.3em,labelsep=0.6em,itemsep=6pt}
% Checklist: quadrado grande para marcar à mão.
\DeclareRobustCommand{\caixinha}{\tikz[baseline=-0.6ex]\draw[cor,line width=1.2pt,rounded corners=2pt](0,-0.16)rectangle(0.42,0.26);}
\newlist{checklist}{itemize}{1}
\setlist[checklist]{label=\caixinha,leftmargin=2em,labelsep=0.6em,itemsep=7pt}

\pagestyle{fancy}
\fancyhf{}
\renewcommand{\headrulewidth}{0pt}
\fancyfoot[L]{\footnotesize\color{cor}\textbf{\faChartLine\ \ Frente C}\color{cinza}\ \ ·\ \ Metas e Indicadores · Guia do projeto}
\fancyfoot[R]{\footnotesize\color{cinza}\textbf{\thepage}\,/\,\pageref{LastPage}}
% Fita da cor da equipe na margem esquerda de todas as páginas.
\fancyhead[L]{\begin{tikzpicture}[remember picture,overlay]\fill[cor](current page.north west)rectangle([xshift=5pt]current page.south west);\end{tikzpicture}}

% Título de seção com ícone: nunca sozinho no pé da página.
\newcommand{\secao}[2]{\Needspace{10\baselineskip}\vspace{12pt}%
  \noindent\tikz[baseline=(i.base)]\node[circle,fill=cor,text=white,minimum size=26pt,inner sep=0pt](i){\small #1};%
  \hspace{8pt}{\Large\bfseries #2}\par\vspace{-4pt}{\color{cor!35}\rule{\linewidth}{1.2pt}}\par\vspace{2pt}}
\newcommand{\subsecao}[1]{\Needspace{8\baselineskip}\vspace{8pt}{\large\bfseries\color{cor}#1}\par}
\newcommand{\rotulo}[2]{\Needspace{6\baselineskip}\vspace{10pt}%
  \noindent{\color{cor}\normalsize #1}\hspace{6pt}{\footnotesize\bfseries\color{cor}\MakeUppercase{#2}}\par\vspace{2pt}}
\newtcolorbox{destaque}[1][]{enhanced,breakable,colback=suave,colframe=suave,borderline west={4pt}{0pt}{cor},arc=4pt,boxrule=0pt,left=12pt,right=12pt,top=8pt,bottom=8pt,before skip=8pt,after skip=10pt,#1}
\newtcolorbox{apostila}{enhanced,breakable,colback=white,colframe=cor,boxrule=1pt,arc=6pt,left=12pt,right=12pt,top=8pt,bottom=8pt,before skip=8pt,after skip=8pt,
  overlay={\node[fill=cor,text=white,rounded corners=3pt,font=\scriptsize\bfseries,inner sep=4pt,anchor=west] at ([xshift=10pt]frame.north west) {\faBookOpen\ DA APOSTILA};}}
\newtcolorbox{selo}{enhanced,left=8pt,right=8pt,top=8pt,bottom=8pt,before skip=0pt,after skip=0pt,
  frame code={\draw[cor,dashed,line width=1.4pt,rounded corners=8pt](frame.south west)rectangle(frame.north east);},
  interior code={\fill[suave,rounded corners=8pt](interior.south west)rectangle(interior.north east);}}
% Anotações: ocupa o resto da página da etapa (só se sobrar espaço útil).
\newcommand{\anotacoes}{\par\ifdim\dimexpr\pagegoal-\pagetotal\relax>3.5cm\ifdim\pagegoal<\maxdimen
\begin{tcolorbox}[enhanced,height fill,colback=white,colframe=cor!35,boxrule=0.8pt,arc=8pt,before skip=10pt,left=10pt,top=6pt,
  borderline={0.8pt}{0pt}{cor!35,dashed},frame hidden,
  underlay={\begin{scope}\clip(interior.south west)rectangle(interior.north east);
    \foreach \y in {1,...,40}{\draw[cor!18,line width=0.6pt]([yshift=-\y*0.75cm-0.5cm]interior.north west)--([yshift=-\y*0.75cm-0.5cm]interior.north east);}\end{scope}}]
{\footnotesize\bfseries\color{cor}\faStickyNote\ \ ANOTAÇÕES}
\end{tcolorbox}\fi\fi}
% Fim da etapa: se o bloco final não cabe, a sobra vira Anotações e o bloco vai para a próxima página.
\newif\ifjaanotou
\newsavebox{\blocofinal}
% Mede o bloco final já montado; se não couber no resto da página, a sobra vira Anotações.
\newcommand{\usarblocofinal}{\par\global\jaanotoufalse
  \ifdim\pagegoal<\maxdimen\ifdim\dimexpr\pagegoal-\pagetotal\relax<\dimexpr\ht\blocofinal+\dp\blocofinal+18pt\relax
    \anotacoes\global\jaanotoutrue\clearpage\fi\fi
  \vspace{10pt}\noindent\usebox{\blocofinal}\par}
\newtcolorbox{faixa}[1][]{enhanced,colback=cor,colframe=cor,coltext=white,arc=8pt,boxrule=0pt,left=14pt,right=14pt,top=10pt,bottom=10pt,#1}
\newtcbox{\chip}[1][]{enhanced,on line,colback=suave,colframe=suave,coltext=cor,arc=7pt,boxrule=0pt,left=5pt,right=5pt,top=2pt,bottom=2pt,fontupper=\footnotesize\bfseries,#1}
\newtcbox{\chipbranco}{enhanced,on line,colback=white,colframe=white,coltext=cor,arc=7pt,boxrule=0pt,left=6pt,right=6pt,top=2pt,bottom=2pt,fontupper=\footnotesize\bfseries}
\begin{document}
\thispagestyle{fancy}
\begin{tikzpicture}[remember picture,overlay]
\fill[cor,opacity=0.10] ([xshift=-1.5cm,yshift=-1cm]current page.north east) circle (5.2cm);
\fill[cor,opacity=0.07] ([xshift=-5.5cm,yshift=-4.5cm]current page.north east) circle (2.4cm);
\fill[cor,opacity=0.06] ([xshift=2cm,yshift=3cm]current page.south west) circle (4cm);
\end{tikzpicture}
\vspace*{2mm}
\noindent\begin{tikzpicture}
\node[circle,fill=cor,text=white,minimum size=78pt,font=\fontsize{40}{40}\selectfont\bfseries] (L) {C};
\node[right=14pt of L.east,anchor=west,align=left,text width=12.2cm] {
  {\footnotesize\bfseries\color{cor}PROJETO ESTADIA\ \ ·\ \ EQUIPE DOURADO\ \ ·\ \ FRENTE C}\\[4pt]
  {\fontsize{23}{27}\selectfont\bfseries Metas e Indicadores}\\[6pt]
  \chip{\faCompass\ \ Guia do projeto}\ \chip{\faChartLine\ \ Aulas 10, 11 e 12 da apostila}
};
\end{tikzpicture}
\vspace{10pt}

\begin{faixa}
{\large\faRocket}\hspace{8pt}{\large\bfseries Nossa missão}\par\vspace{2pt}
{\normalsize Mostrar com números como o Estadia ajuda o hotel: definir metas e os indicadores que a dona vai acompanhar.}
\end{faixa}

\vspace{6pt}
\begin{tcolorbox}[enhanced,colback=white,colframe=cor!40,boxrule=1pt,arc=8pt,left=10pt,right=10pt,top=8pt,bottom=10pt,
  title={\faUserFriends\ \ Nossa equipe},fonttitle=\bfseries\small,coltitle=cor,colbacktitle=suave,attach boxed title to top left={xshift=10pt,yshift=-8pt},boxed title style={arc=4pt,boxrule=0pt}]
\vspace{6pt}
\begin{minipage}[t]{0.31\linewidth}\centering
\tikz\node[circle,draw=cor,line width=1.4pt,fill=white,minimum size=34pt,text=cor]{\large \faBullseye};\par\vspace{4pt}
{\scriptsize\bfseries\color{cor}\MakeUppercase{Metas}}\par\vspace{14pt}
{\color{cor!60}\rule{0.9\linewidth}{0.8pt}}\par{\scriptsize\color{cinza}nome}
\end{minipage}\hfill
\begin{minipage}[t]{0.31\linewidth}\centering
\tikz\node[circle,draw=cor,line width=1.4pt,fill=white,minimum size=34pt,text=cor]{\large \faChartBar};\par\vspace{4pt}
{\scriptsize\bfseries\color{cor}\MakeUppercase{Painel}}\par\vspace{14pt}
{\color{cor!60}\rule{0.9\linewidth}{0.8pt}}\par{\scriptsize\color{cinza}nome}
\end{minipage}\hfill
\begin{minipage}[t]{0.31\linewidth}\centering
\tikz\node[circle,draw=cor,line width=1.4pt,fill=white,minimum size=34pt,text=cor]{\large \faLaptopCode};\par\vspace{4pt}
{\scriptsize\bfseries\color{cor}\MakeUppercase{Dev de apoio}}\par\vspace{14pt}
{\color{cor!60}\rule{0.9\linewidth}{0.8pt}}\par{\scriptsize\color{cinza}nome}
\end{minipage}
\end{tcolorbox}
\begin{destaque}[before skip=4pt]
{\color{cor}\Large\faCrosshairs}\hspace{8pt}{\large\bfseries O seu foco}\par\vspace{3pt}
Vocês transformam o sistema em números que ajudam a dona do hotel a decidir: subir ou baixar preço, fazer promoção, contratar.
\end{destaque}
\secao{\faSitemap}{Como usar este guia}
{\large Este guia tem uma parte geral e um \textbf{manual para cada função}. Leia a parte geral e depois vá direto ao manual da sua função: ele tem tudo de que você precisa para fazer as suas tarefas do Passo a passo.}\par\vspace{6pt}
\Needspace{7\baselineskip}
\begin{tcbraster}[raster columns=3,raster equal height=rows,raster column skip=6pt,raster row skip=6pt,raster before skip=4pt,raster after skip=8pt]
\begin{tcolorbox}[enhanced,colback=white,colframe=cor!25,boxrule=0.8pt,arc=6pt,left=8pt,right=8pt,top=6pt,bottom=6pt,colback=suave!40]
\centering{\color{cor}\LARGE \faBullseye}\par\vspace{3pt}{\bfseries Metas}\par\vspace{2pt}{\footnotesize\color{cinza}Define as metas do hotel e explica os indicadores.}\par\vspace{4pt}\chip{C1}\ \chip{C2}\par\vspace{4pt}{\footnotesize\bfseries\color{cor}\faArrowRight\ manual na página \pageref{manual-c1}}
\end{tcolorbox}
\begin{tcolorbox}[enhanced,colback=white,colframe=cor!25,boxrule=0.8pt,arc=6pt,left=8pt,right=8pt,top=6pt,bottom=6pt,colback=suave!40]
\centering{\color{cor}\LARGE \faChartBar}\par\vspace{3pt}{\bfseries Painel}\par\vspace{2pt}{\footnotesize\color{cinza}Monta a planilha com os números, as cores e o gráfico.}\par\vspace{4pt}\chip{C3}\par\vspace{4pt}{\footnotesize\bfseries\color{cor}\faArrowRight\ manual na página \pageref{manual-c2}}
\end{tcolorbox}
\begin{tcolorbox}[enhanced,colback=white,colframe=cor!25,boxrule=0.8pt,arc=6pt,left=8pt,right=8pt,top=6pt,bottom=6pt,colback=suave!40]
\centering{\color{cor}\LARGE \faLaptopCode}\par\vspace{3pt}{\bfseries Dev de apoio}\par\vspace{2pt}{\footnotesize\color{cinza}Passa os números da versão de demonstração.}\par\vspace{4pt}\chip{apoio às tarefas}\par\vspace{4pt}{\footnotesize\bfseries\color{cor}\faArrowRight\ manual na página \pageref{manual-c-dev}}
\end{tcolorbox}
\end{tcbraster}
\secao{\faHotel}{O projeto em 1 minuto}
{\large O Estadia é um sistema para administrar um hotel de 50 quartos. Pelo computador ou pelo celular, a equipe do hotel faz reservas, registra a entrada (check-in) e a saída (check-out) dos hóspedes, cobra o que eles consumiram, fecha o caixa do dia e vê relatórios.}\par
\Needspace{7\baselineskip}
\begin{tcbraster}[raster columns=2,raster equal height=rows,raster column skip=6pt,raster row skip=6pt,raster before skip=4pt,raster after skip=8pt]
\begin{tcolorbox}[enhanced,colback=white,colframe=cor!25,boxrule=0.8pt,arc=6pt,left=8pt,right=8pt,top=6pt,bottom=6pt,colback=suave!50]
{\color{cor}\Large \faExclamationTriangle}\hspace{6pt}{\bfseries O problema}\par\vspace{3pt}{\small Muitos hotéis pequenos ainda controlam tudo em caderno, planilha e WhatsApp. Aí acontecem erros: vender o mesmo quarto para duas pessoas, cobrar errado, o caixa não bater e a dona não saber quanto o hotel realmente ganhou no mês.}
\end{tcolorbox}
\begin{tcolorbox}[enhanced,colback=white,colframe=cor!25,boxrule=0.8pt,arc=6pt,left=8pt,right=8pt,top=6pt,bottom=6pt,colback=suave!50]
{\color{cor}\Large \faLightbulb}\hspace{6pt}{\bfseries A solução}\par\vspace{3pt}{\small Um sistema único que vários funcionários usam ao mesmo tempo. Ele não deixa vender o mesmo quarto duas vezes, guarda quem fez cada ação e mostra relatórios prontos para a dona decidir preços e equipe.}
\end{tcolorbox}
\end{tcbraster}
\subsecao{\faAngleDoubleRight\ Os 4 indicadores, com um exemplo resolvido}
Um hotel de 50 quartos, num mês de 30 dias, tem 1.500 ``noites de quarto'' para vender. Imagine que ele vendeu 1.050 noites e recebeu R\$ 252.000 de diárias:\par
\begingroup\small
\begin{longtable}{>{\raggedright\arraybackslash}p{0.193\linewidth}>{\raggedright\arraybackslash}p{0.494\linewidth}>{\raggedright\arraybackslash}p{0.205\linewidth}}
\rowcolor{cor} \color{white}\textbf{\footnotesize Indicador} & \color{white}\textbf{\footnotesize Conta} & \color{white}\textbf{\footnotesize Resultado} \\
\endfirsthead
\rowcolor{cor} \color{white}\textbf{\footnotesize Indicador} & \color{white}\textbf{\footnotesize Conta} & \color{white}\textbf{\footnotesize Resultado} \\
\endhead
\rowcolor{white}\textbf{Taxa de ocupação} & 1.050 ÷ 1.500 & 70\% \\ \hline
\rowcolor{suave!55}\textbf{Diária média} & R\$ 252.000 ÷ 1.050 & R\$ 240,00 \\ \hline
\rowcolor{white}\textbf{RevPAR} & R\$ 252.000 ÷ 1.500 & R\$ 168,00 \\ \hline
\rowcolor{suave!55}\textbf{No-show} & reservas que não vieram ÷ reservas do mês & ex.: 6 ÷ 120 = 5\% \\ \hline
\end{longtable}
\endgroup
\subsecao{\faAngleDoubleRight\ Onde esses números aparecem no sistema}
\begin{destaque}\begin{itemize}
  \item Tela Relatórios → ``Ocupação, diária média e RevPAR''.
  \item Relatório de no-shows (quem reservou e não veio).
  \item Todos os relatórios podem ser baixados em planilha (CSV).
\end{itemize}
\end{destaque}
\clearpage
\begin{faixa}[before skip=0pt]
\begin{tabular}{@{}m{46pt}>{\raggedright\arraybackslash}m{0.76\linewidth}@{}}
\tikz\node[circle,fill=white,text=cor,minimum size=42pt,inner sep=0pt]{\LARGE \faBullseye}; &
{\footnotesize\bfseries MANUAL DA FUNÇÃO · FRENTE C}\label{manual-c1}\par\vspace{2pt}
{\fontsize{20}{24}\selectfont\bfseries Metas}\par\vspace{3pt}
{\small Define as metas do hotel e explica os indicadores.}
\end{tabular}
\end{faixa}
\rotulo{\faTasks}{Suas tarefas}
\chip{C1}\ Metas do hotel (OKR)\hspace{10pt}\chip{C2}\ Fichas de indicadores (KPIs)\par
{\footnotesize\color{cinza}Você também ajuda nas outras tarefas da frente. Quem faz cada uma está no Passo a passo.}\par
\Needspace{12\baselineskip}\vspace{10pt}
\begin{tcolorbox}[enhanced,colback=cor,colframe=cor,coltext=white,arc=6pt,boxrule=0pt,left=10pt,right=10pt,top=6pt,bottom=6pt,after skip=6pt]
{\Large \faBullseye}\hspace{10pt}{\footnotesize\bfseries TAREFA C1}\hspace{8pt}{\large\bfseries Metas do hotel (OKR)}\hfill{\footnotesize\faClock\ cerca de 1h}
\end{tcolorbox}
\rotulo{\faBrain}{O que você precisa saber}
\begin{itemize}
  \item Objetivo: frase inspiradora, SEM números.
  \item Resultado-chave: número de antes → número da meta. Não é tarefa.
\end{itemize}
\Needspace{7\baselineskip}
\begin{tcbraster}[raster columns=2,raster equal height=rows,raster column skip=6pt,raster row skip=6pt,raster before skip=4pt,raster after skip=8pt]
\begin{tcolorbox}[enhanced,colback=white,colframe=cor!25,boxrule=0.8pt,arc=6pt,left=8pt,right=8pt,top=6pt,bottom=6pt,colback=suave!50]
{\footnotesize\bfseries\color{cor}\MakeUppercase{Dicas}}\par\vspace{3pt}{\small {\color{cor}\faLightbulb}\ \ ``Antes'' pode ser estimado: escrevam a fonte (palpite da dona, observação, pesquisa).}
\end{tcolorbox}
\begin{tcolorbox}[enhanced,colback=white,colframe=cor!25,boxrule=0.8pt,arc=6pt,left=8pt,right=8pt,top=6pt,bottom=6pt,colback=suave!50]
{\footnotesize\bfseries\color{cor}\MakeUppercase{Cuidado com}}\par\vspace{3pt}{\small {\color{cor}\faExclamationTriangle}\ \ Resultado-chave como tarefa: ``instalar o sistema'' é tarefa; ``check-in em 3 minutos'' é resultado.}
\end{tcolorbox}
\end{tcbraster}
\subsecao{\faBoxOpen\ Ideias de resultados-chave para o hotel (escolham 3)}
\begingroup\small
\begin{longtable}{>{\raggedright\arraybackslash}p{0.532\linewidth}>{\raggedright\arraybackslash}p{0.250\linewidth}>{\raggedright\arraybackslash}p{0.110\linewidth}}
\rowcolor{cor} \color{white}\textbf{\footnotesize Resultado-chave} & \color{white}\textbf{\footnotesize Antes (estimado)} & \color{white}\textbf{\footnotesize Meta} \\
\endfirsthead
\rowcolor{cor} \color{white}\textbf{\footnotesize Resultado-chave} & \color{white}\textbf{\footnotesize Antes (estimado)} & \color{white}\textbf{\footnotesize Meta} \\
\endhead
\rowcolor{white}\textbf{Tempo médio de check-in} & 8 min & 3 min \\ \hline
\rowcolor{suave!55}\textbf{Reservas vendidas em dobro por mês} & 3 & 0 \\ \hline
\rowcolor{white}\textbf{Erros de cobrança por mês} & 10 & 2 \\ \hline
\rowcolor{suave!55}\textbf{Tempo para fechar o caixa} & 40 min & 15 min \\ \hline
\rowcolor{white}\textbf{Dias com o fechamento do dia feito} & 50\% & 100\% \\ \hline
\end{longtable}
\endgroup
\subsecao{\faBoxOpen\ Exemplos de objetivo}
\begin{destaque}\begin{itemize}
  \item Deixar a recepção rápida e sem erros.
  \item Dar à dona do hotel controle total do dia a dia.
  \item Acabar com o caderno e a planilha na recepção.
\end{itemize}
\end{destaque}
\Needspace{12\baselineskip}\vspace{10pt}
\begin{tcolorbox}[enhanced,colback=cor,colframe=cor,coltext=white,arc=6pt,boxrule=0pt,left=10pt,right=10pt,top=6pt,bottom=6pt,after skip=6pt]
{\Large \faStopwatch}\hspace{10pt}{\footnotesize\bfseries TAREFA C2}\hspace{8pt}{\large\bfseries Fichas de indicadores (KPIs)}\hfill{\footnotesize\faClock\ cerca de 1h30}
\end{tcolorbox}
\rotulo{\faBrain}{O que você precisa saber}
\begin{itemize}
  \item Um bom indicador ajuda a decidir alguma coisa.
  \item Os 4 indicadores já são calculados pelo sistema (tela Relatórios).
\end{itemize}
\Needspace{7\baselineskip}
\begin{tcbraster}[raster columns=2,raster equal height=rows,raster column skip=6pt,raster row skip=6pt,raster before skip=4pt,raster after skip=8pt]
\begin{tcolorbox}[enhanced,colback=white,colframe=cor!25,boxrule=0.8pt,arc=6pt,left=8pt,right=8pt,top=6pt,bottom=6pt,colback=suave!50]
{\footnotesize\bfseries\color{cor}\MakeUppercase{Dicas}}\par\vspace{3pt}{\small {\color{cor}\faLightbulb}\ \ Refaçam a conta de exemplo na calculadora antes de preencher: se bater, entenderam.}
\end{tcolorbox}
\begin{tcolorbox}[enhanced,colback=white,colframe=cor!25,boxrule=0.8pt,arc=6pt,left=8pt,right=8pt,top=6pt,bottom=6pt,colback=suave!50]
{\footnotesize\bfseries\color{cor}\MakeUppercase{Cuidado com}}\par\vspace{3pt}{\small {\color{cor}\faExclamationTriangle}\ \ Confundir diária média (divide pelos quartos vendidos) com RevPAR (divide por todos os quartos).}
\end{tcolorbox}
\end{tcbraster}
\subsecao{\faBoxOpen\ Os 4 indicadores explicados}
\begingroup\small
\begin{longtable}{>{\raggedright\arraybackslash}p{0.121\linewidth}>{\raggedright\arraybackslash}p{0.288\linewidth}>{\raggedright\arraybackslash}p{0.098\linewidth}>{\raggedright\arraybackslash}p{0.348\linewidth}}
\rowcolor{cor} \color{white}\textbf{\footnotesize Indicador} & \color{white}\textbf{\footnotesize Fórmula} & \color{white}\textbf{\footnotesize Meta sugerida} & \color{white}\textbf{\footnotesize Decisão que ajuda} \\
\endfirsthead
\rowcolor{cor} \color{white}\textbf{\footnotesize Indicador} & \color{white}\textbf{\footnotesize Fórmula} & \color{white}\textbf{\footnotesize Meta sugerida} & \color{white}\textbf{\footnotesize Decisão que ajuda} \\
\endhead
\rowcolor{white}\textbf{Taxa de ocupação} & quartos ocupados ÷ quartos do hotel & 70\% & Baixa? Fazer promoção. \\ \hline
\rowcolor{suave!55}\textbf{Diária média} & receita das diárias ÷ quartos vendidos & R\$ 240 & Muito baixa? Rever descontos. \\ \hline
\rowcolor{white}\textbf{RevPAR} & receita das diárias ÷ todos os quartos & R\$ 168 & Mostra se preço e ocupação estão equilibrados. \\ \hline
\rowcolor{suave!55}\textbf{No-show} & reservas que não vieram ÷ reservas & até 5\% & Alto? Cobrar sinal na reserva. \\ \hline
\end{longtable}
\endgroup
\subsecao{\faBoxOpen\ Conta de exemplo (um mês)}
\begin{destaque}{\color{cor}\faQuoteLeft}\ \ 50 quartos × 30 dias = 1.500 noites. Vendidas: 1.050 noites. Receita das diárias: R\$ 252.000. Ocupação = 1.050 ÷ 1.500 = 70\%. Diária média = 252.000 ÷ 1.050 = R\$ 240. RevPAR = 252.000 ÷ 1.500 = R\$ 168.\end{destaque}
\clearpage
\begin{faixa}[before skip=0pt]
\begin{tabular}{@{}m{46pt}>{\raggedright\arraybackslash}m{0.76\linewidth}@{}}
\tikz\node[circle,fill=white,text=cor,minimum size=42pt,inner sep=0pt]{\LARGE \faChartBar}; &
{\footnotesize\bfseries MANUAL DA FUNÇÃO · FRENTE C}\label{manual-c2}\par\vspace{2pt}
{\fontsize{20}{24}\selectfont\bfseries Painel}\par\vspace{3pt}
{\small Monta a planilha com os números, as cores e o gráfico.}
\end{tabular}
\end{faixa}
\rotulo{\faTasks}{Suas tarefas}
\chip{C3}\ Painel de indicadores\par
{\footnotesize\color{cinza}Você também ajuda nas outras tarefas da frente. Quem faz cada uma está no Passo a passo.}\par
\Needspace{12\baselineskip}\vspace{10pt}
\begin{tcolorbox}[enhanced,colback=cor,colframe=cor,coltext=white,arc=6pt,boxrule=0pt,left=10pt,right=10pt,top=6pt,bottom=6pt,after skip=6pt]
{\Large \faChartBar}\hspace{10pt}{\footnotesize\bfseries TAREFA C3}\hspace{8pt}{\large\bfseries Painel de indicadores}\hfill{\footnotesize\faClock\ cerca de 2h}
\end{tcolorbox}
\rotulo{\faBrain}{O que você precisa saber}
\begin{itemize}
  \item Painel bom se entende em 30 segundos: número, meta e cor.
\end{itemize}
\Needspace{7\baselineskip}
\begin{tcbraster}[raster columns=2,raster equal height=rows,raster column skip=6pt,raster row skip=6pt,raster before skip=4pt,raster after skip=8pt]
\begin{tcolorbox}[enhanced,colback=white,colframe=cor!25,boxrule=0.8pt,arc=6pt,left=8pt,right=8pt,top=6pt,bottom=6pt,colback=suave!50]
{\footnotesize\bfseries\color{cor}\MakeUppercase{Dicas}}\par\vspace{3pt}{\small {\color{cor}\faLightbulb}\ \ Ponham ícones ou a palavra (OK / Atenção / Problema) junto da cor: nem todo mundo enxerga cores igual.}
\end{tcolorbox}
\begin{tcolorbox}[enhanced,colback=white,colframe=cor!25,boxrule=0.8pt,arc=6pt,left=8pt,right=8pt,top=6pt,bottom=6pt,colback=suave!50]
{\footnotesize\bfseries\color{cor}\MakeUppercase{Cuidado com}}\par\vspace{3pt}{\small {\color{cor}\faExclamationTriangle}\ \ Gráfico sem título e sem unidade (\%). Sempre digam o que é o número.}
\end{tcolorbox}
\end{tcbraster}
\subsecao{\faBoxOpen\ Números de exemplo para a planilha (4 semanas)}
\begingroup\small
\begin{longtable}{>{\raggedright\arraybackslash}p{0.176\linewidth}>{\raggedright\arraybackslash}p{0.201\linewidth}>{\raggedright\arraybackslash}p{0.302\linewidth}>{\raggedright\arraybackslash}p{0.176\linewidth}}
\rowcolor{cor} \color{white}\textbf{\footnotesize Semana} & \color{white}\textbf{\footnotesize Ocupação} & \color{white}\textbf{\footnotesize Diária média} & \color{white}\textbf{\footnotesize No-show} \\
\endfirsthead
\rowcolor{cor} \color{white}\textbf{\footnotesize Semana} & \color{white}\textbf{\footnotesize Ocupação} & \color{white}\textbf{\footnotesize Diária média} & \color{white}\textbf{\footnotesize No-show} \\
\endhead
\rowcolor{white}\textbf{1} & 62\% & R\$ 230 & 6\% \\ \hline
\rowcolor{suave!55}\textbf{2} & 68\% & R\$ 236 & 5\% \\ \hline
\rowcolor{white}\textbf{3} & 74\% & R\$ 244 & 4\% \\ \hline
\rowcolor{suave!55}\textbf{4} & 76\% & R\$ 248 & 4\% \\ \hline
\end{longtable}
\endgroup
\subsecao{\faBoxOpen\ Google Sheets passo a passo}
\begin{destaque}\begin{enumerate}
  \item Digitem a tabela acima.
  \item Selecionem as colunas Semana e Ocupação → Inserir → Gráfico (linha).
  \item Formatar → Formatação condicional: verde se atingiu a meta, amarelo se está perto, vermelho se está longe.
  \item No no-show é ao contrário: quanto menor, melhor.
\end{enumerate}
\end{destaque}
\subsecao{\faBoxOpen\ Modelo de análise (3 frases)}
\begin{destaque}{\color{cor}\faQuoteLeft}\ \ A ocupação subiu de 62\% para 76\% em 4 semanas e passou da meta. O no-show caiu para 4\%. Recomendamos manter os preços e cobrar sinal nas reservas de fim de semana.\end{destaque}
\Needspace{22\baselineskip}\vspace{14pt}
\begin{faixa}[before skip=0pt]
\begin{tabular}{@{}m{46pt}>{\raggedright\arraybackslash}m{0.76\linewidth}@{}}
\tikz\node[circle,fill=white,text=cor,minimum size=42pt,inner sep=0pt]{\LARGE \faLaptopCode}; &
{\footnotesize\bfseries MANUAL DA FUNÇÃO · FRENTE C}\label{manual-c-dev}\par\vspace{2pt}
{\fontsize{20}{24}\selectfont\bfseries Dev de apoio}\par\vspace{3pt}
{\small Passa os números da versão de demonstração.}
\end{tabular}
\end{faixa}
\secao{\faTools}{O que você entrega em cada tarefa}
\begingroup\small
\begin{longtable}{>{\raggedright\arraybackslash}p{0.105\linewidth}>{\raggedright\arraybackslash}p{0.823\linewidth}}
\rowcolor{cor} \color{white}\textbf{\footnotesize Tarefa} & \color{white}\textbf{\footnotesize O que fazer} \\
\endfirsthead
\rowcolor{cor} \color{white}\textbf{\footnotesize Tarefa} & \color{white}\textbf{\footnotesize O que fazer} \\
\endhead
\rowcolor{white}\textbf{C1} & Confirmar que o sistema consegue medir os resultados-chave escolhidos. \\ \hline
\rowcolor{suave!55}\textbf{C2} & Mostrar a tela Relatórios → ``Ocupação, diária média e RevPAR''. \\ \hline
\rowcolor{white}\textbf{C3} & Exportar um relatório em CSV da versão de demonstração (ou confirmar os números de exemplo). \\ \hline
\end{longtable}
\endgroup
\rotulo{\faHandsHelping}{Sempre}
\begin{itemize}
  \item Mostrar o sistema funcionando e criar usuários de teste.
  \item Conferir o que a frente escreveu sobre o sistema.
  \item Passar números e prints da versão de demonstração.
  \item Explicar o que for técnico em palavras simples.
\end{itemize}
\begin{destaque}{\color{cor}\faDesktop}\ \ Peça ao seu dev de apoio o endereço da versão de demonstração e um usuário de teste. Ela tem dados inventados: nunca use dados reais de hóspedes nos trabalhos ou prints.\end{destaque}
\clearpage
\secao{\faQuestionCircle}{Onde tirar dúvidas}
\begingroup\small
\begin{longtable}{>{\raggedright\arraybackslash}p{0.342\linewidth}>{\raggedright\arraybackslash}p{0.586\linewidth}}
\rowcolor{cor} \color{white}\textbf{\footnotesize Preciso de…} & \color{white}\textbf{\footnotesize Onde encontrar} \\
\endfirsthead
\rowcolor{cor} \color{white}\textbf{\footnotesize Preciso de…} & \color{white}\textbf{\footnotesize Onde encontrar} \\
\endhead
\rowcolor{white}\textbf{Ver o sistema funcionando} & Peça ao dev de apoio o link da versão de demonstração \\ \hline
\rowcolor{suave!55}\textbf{Tirar dúvida sobre o sistema} & Dev de apoio da sua frente \\ \hline
\rowcolor{white}\textbf{Conceitos da disciplina} & Apostila ``Gestão Estratégica Guiada a Dados'' \\ \hline
\rowcolor{suave!55}\textbf{O que a dona do hotel respondeu} & Frente D (questionário da reunião) \\ \hline
\rowcolor{white}\textbf{Entregar seu trabalho} & Site da divisão de funções → Entregas da sua frente \\ \hline
\end{longtable}
\endgroup
\secao{\faBook}{Glossário}
\Needspace{27\baselineskip}
\begin{tcbraster}[raster columns=2,raster equal height=rows,raster column skip=6pt,raster row skip=6pt,raster before skip=4pt,raster after skip=8pt]
\begin{tcolorbox}[enhanced,colback=white,colframe=cor!25,boxrule=0.8pt,arc=6pt,left=8pt,right=8pt,top=6pt,bottom=6pt]
{\bfseries\color{cor}Diária}\par\vspace{1pt}{\footnotesize O preço de uma noite no hotel.}
\end{tcolorbox}
\begin{tcolorbox}[enhanced,colback=white,colframe=cor!25,boxrule=0.8pt,arc=6pt,left=8pt,right=8pt,top=6pt,bottom=6pt]
{\bfseries\color{cor}Check-in / Check-out}\par\vspace{1pt}{\footnotesize Registro da entrada e da saída do hóspede.}
\end{tcolorbox}
\begin{tcolorbox}[enhanced,colback=white,colframe=cor!25,boxrule=0.8pt,arc=6pt,left=8pt,right=8pt,top=6pt,bottom=6pt]
{\bfseries\color{cor}Taxa de ocupação}\par\vspace{1pt}{\footnotesize Quantos \% dos quartos foram ocupados. Ex.: 35 de 50 = 70\%.}
\end{tcolorbox}
\begin{tcolorbox}[enhanced,colback=white,colframe=cor!25,boxrule=0.8pt,arc=6pt,left=8pt,right=8pt,top=6pt,bottom=6pt]
{\bfseries\color{cor}Diária média (ADR)}\par\vspace{1pt}{\footnotesize Quanto o hotel recebeu, em média, por quarto vendido.}
\end{tcolorbox}
\begin{tcolorbox}[enhanced,colback=white,colframe=cor!25,boxrule=0.8pt,arc=6pt,left=8pt,right=8pt,top=6pt,bottom=6pt]
{\bfseries\color{cor}RevPAR}\par\vspace{1pt}{\footnotesize Receita das diárias dividida por todos os quartos (vendidos ou não).}
\end{tcolorbox}
\begin{tcolorbox}[enhanced,colback=white,colframe=cor!25,boxrule=0.8pt,arc=6pt,left=8pt,right=8pt,top=6pt,bottom=6pt]
{\bfseries\color{cor}No-show}\par\vspace{1pt}{\footnotesize Quem reservou e não apareceu.}
\end{tcolorbox}
\begin{tcolorbox}[enhanced,colback=white,colframe=cor!25,boxrule=0.8pt,arc=6pt,left=8pt,right=8pt,top=6pt,bottom=6pt]
{\bfseries\color{cor}Governança}\par\vspace{1pt}{\footnotesize Equipe de limpeza e arrumação dos quartos.}
\end{tcolorbox}
\begin{tcolorbox}[enhanced,colback=white,colframe=cor!25,boxrule=0.8pt,arc=6pt,left=8pt,right=8pt,top=6pt,bottom=6pt]
{\bfseries\color{cor}Backlog}\par\vspace{1pt}{\footnotesize Lista de tudo que o sistema precisa ter, em ordem de importância.}
\end{tcolorbox}
\begin{tcolorbox}[enhanced,colback=white,colframe=cor!25,boxrule=0.8pt,arc=6pt,left=8pt,right=8pt,top=6pt,bottom=6pt]
{\bfseries\color{cor}Sprint}\par\vspace{1pt}{\footnotesize Um ciclo curto de trabalho (aqui, 1 semana) com um objetivo.}
\end{tcolorbox}
\begin{tcolorbox}[enhanced,colback=white,colframe=cor!25,boxrule=0.8pt,arc=6pt,left=8pt,right=8pt,top=6pt,bottom=6pt]
{\bfseries\color{cor}KPI}\par\vspace{1pt}{\footnotesize Indicador importante, que ajuda a tomar uma decisão.}
\end{tcolorbox}
\end{tcbraster}
\end{document}
